% ============================================================================
% CAPITOLO 5 — RISULTATI SPERIMENTALI E VALIDAZIONE
% ============================================================================
\chapter{Risultati e validazione}
\label{cap:risultati}

% ----------------------------------------------------------------------------
% NOTA METODOLOGICA FONDAMENTALE SULLA STESURA
% ----------------------------------------------------------------------------
\noindent\fbox{
\parbox{0.96\textwidth}{
\textbf{\large REGOLE DI STESURA PER LO STUDENTE — LEGGERE PRIMA DI SCRIVERE:}
\vspace{0.2cm}
\begin{enumerate}
    \item \textbf{SCRIVERE QUESTO CAPITOLO SUBITO DOPO IL CAPITOLO 4:} Non appena hai raccolto i dati dei test e generato i grafici, \textbf{redigi immediatamente questo capitolo}!
    \item \textbf{PERCHÉ?} I numeri, le configurazioni sperimentali, i parametri dei test e il significato delle anomalie osservate sono ancora chiari e verificabili. Se aspetti, rischi di dimenticare le condizioni esatte in cui sono stati eseguiti i test.
    \item \textbf{STRUTTURA SCIENTIFICA:} Un capitolo di risultati deve sempre seguire questo schema logico:
    \begin{itemize}
        \item Prima si definisce la \textbf{metodologia di test};
        \item Poi si definiscono formalmente le \textbf{metriche matematiche};
        \item Poi si mostrano i \textbf{dati quantitativi} (tabelle e grafici);
        \item Infine si conduce una \textbf{discussione critica} approfondita (senza limitarsi a descrivere i grafici).
    \end{itemize}
\end{enumerate}
}
}
\vspace{0.6cm}

Questo capitolo presenta la metodologia di validazione adottata e i risultati sperimentali ottenuti per verificare l'efficacia, la correttezza e la robustezza del sistema proposto.


% ============================================================================
\section{Metodologia di Validazione e Setup Sperimentale}
\label{sec:metodologia-validazione}
% ============================================================================

\textbf{Guida per lo studente:} Descrivere dettagliatamente l'ambiente e gli scenari in cui sono stati svolti i test per garantire la riproducibilità scientifica:
\begin{itemize}
    \item \textbf{Ambiente hardware e software di test:} Specifiche della macchina (CPU, RAM, OS, versione runtime);
    \item \textbf{Dataset o scenari di simulazione:} Descrizione dei dati utilizzati (dimensioni, provenienza, split train/test o parametri di generazione sintetica);
    \item \textbf{Scenari di test strutturati:}
    \begin{enumerate}
        \item \emph{Funzionamento nominale (Baseline):} Verifica del comportamento del sistema in condizioni ordinarie prive di disturbi;
        \item \emph{Analisi dei componenti (Ablation Study):} Rimozione selettiva di singoli moduli per verificare se ciascuno di essi è strettamente necessario;
        \item \emph{Stress Test e casi limite:} Valutazione del comportamento del sistema a fronte di input malformati, rumore estremo o attacchi malevoli coordinati.
    \end{enumerate}
\end{itemize}


% ============================================================================
\section{Definizione delle Metriche di Valutazione}
\label{sec:metriche-valutazione}
% ============================================================================

\textbf{Guida per lo studente:} Tutte le metriche utilizzate nelle tabelle e nei grafici devono essere definite formalmente \textbf{prima} di presentare i risultati numerici~\cite{powers2020evaluation}.

Nel caso di sistemi di rilevamento e classificazione, le metriche standard basate sulla matrice di confusione sono definite come segue:

\begin{itemize}
    \item \textbf{Precision (Precisione):} Frazione di allarmi generati che corrispondono a reali anomalie:
    \begin{equation}
    \mathrm{Precision} = \frac{\mathrm{TP}}{\mathrm{TP} + \mathrm{FP}}
    \label{eq:precision}
    \end{equation}
    
    \item \textbf{Recall (Sensibilità):} Capacità del sistema di intercettare tutte le anomalie effettivamente presenti:
    \begin{equation}
    \mathrm{Recall} = \frac{\mathrm{TP}}{\mathrm{TP} + \mathrm{FN}}
    \label{eq:recall}
    \end{equation}
    
    \item \textbf{F1-Score:} Media armonica tra Precision e Recall, utile in presenza di classi fortemente sbilanciate:
    \begin{equation}
    \mathrm{F1} = 2 \cdot \frac{\mathrm{Precision} \cdot \mathrm{Recall}}{\mathrm{Precision} + \mathrm{Recall}}
    \label{eq:f1-score}
    \end{equation}
    
    \item \textbf{Tempo medio di reazione (Latenza $T_{\mathrm{react}}$):} Intervallo temporale intercorso tra l'inizio dell'anomalia e la prima azione cautelativa intrapresa dal sistema.
\end{itemize}


% ============================================================================
\section{Risultati Sperimentali}
\label{sec:risultati-sperimentali}
% ============================================================================

\textbf{Guida per lo studente:} Raccogliere i dati in tabelle riepilogative chiare e commentare l'andamento osservato.

La \Cref{tab:risultati-confronto} riporta il confronto quantitativo delle prestazioni tra la configurazione completa del sistema e le varianti semplificate ottenute tramite ablation study.

\begin{table}[htbp]
\centering
\small
\renewcommand{\arraystretch}{1.3}
\begin{tabular}{l c c c c}
\toprule
\textbf{Configurazione Sistema} & \textbf{Precision} & \textbf{Recall} & \textbf{F1-Score} & \textbf{Latenza (ms)} \\
\midrule
Sistema Completo (Proposto)    & \textbf{1.00}      & \textbf{1.00}   & \textbf{1.00}     & \textbf{12.4} \\
Variante Senza Modello Interno  & 0.42               & 0.65            & 0.51              & 28.1 \\
Variante Senza Termine Epistemico & 0.88             & 0.72            & 0.79              & 15.6 \\
Baseline Tradizionale Black-Box & 0.75               & 0.80            & 0.77              & 24.3 \\
\bottomrule
\end{tabular}
\caption{Confronto delle prestazioni tra la soluzione proposta e le configurazioni di baseline.}
\label{tab:risultati-confronto}
\end{table}

Come si evince dalla \Cref{tab:risultati-confronto}, il sistema completo consegue prestazioni ottimali su tutte le metriche considerate, con un tempo medio di risposta contenuto in circa $12$\,ms.


% ============================================================================
\section{Discussione Critica dei Risultati}
\label{sec:discussione-critica}
% ============================================================================

\textbf{Guida per lo studente:} Questa è la sezione più importante del capitolo. Una tesi di laurea non deve limitarsi a descrivere i dati (``la linea rossa sale e la linea blu scende''), ma deve spiegare il \textbf{perché scientifico} di ciò che si osserva:

\begin{itemize}
    \item \textbf{Interpretazione causale:} Perché una variante ha fallito? Quale componente mancante ha causato i falsi positivi o la degradazione delle prestazioni?
    \item \textbf{Risposta alle ipotesi di partenza:} I risultati confermano le aspettative teoriche formulate all'inizio del lavoro?
    \item \textbf{Correttezza del processo vs correttezza del risultato:} Nel caso di validazione white-box, verificare se il sistema ha preso la decisione corretta per il motivo giusto, oppure se ha ottenuto un risultato positivo per pura casualità (come evidenziato nei test di stress).
    \item \textbf{Trade-off operativi:} Quali compromessi emergono tra efficacia di rilevamento, costo computazionale e reattività del controllore?
\end{itemize}

Sulla base delle evidenze raccolte, il \Cref{cap:conclusioni} traccerà il bilancio finale del lavoro, identificando i limiti dell'approccio e le future direzioni di ricerca.
